%!TEX program = xelatex
\documentclass[11pt, a4paper]{article}

% ─── Packages ────────────────────────────────────────────────────────────────
\usepackage{fontspec}
\usepackage{kotex}
\usepackage[margin=2.5cm]{geometry}
\usepackage{amsmath, amssymb, bm}
\usepackage{graphicx}
\usepackage{booktabs}
\usepackage{multirow}
\usepackage{array}
\usepackage{xcolor}
\usepackage{hyperref}
\usepackage{natbib}
\usepackage{caption}
\usepackage{subcaption}
\usepackage{float}
\usepackage{algorithm}
\usepackage{algpseudocode}
\usepackage{enumitem}
\usepackage{parskip}
\usepackage{titlesec}
\usepackage{setspace}

\setmainhangulfont{Malgun Gothic}
\setsanshangulfont{Malgun Gothic}
\setmonohangulfont{Malgun Gothic}

\onehalfspacing
\setlength{\parindent}{0pt}
\setlength{\parskip}{6pt}

\hypersetup{
    colorlinks=true,
    linkcolor=blue!60!black,
    citecolor=green!50!black,
    urlcolor=blue!70!black
}

\graphicspath{{figures/}}

% ─── Title ───────────────────────────────────────────────────────────────────
\title{
    \vspace{-1cm}
    {\large 캡스톤 디자인 프로젝트 — 최종 보고서}\\[0.8em]
    \textbf{PPO 기반 동적 그리드 트레이딩:\\
    베이지안 계수 최적화와 변동성 적응형 주문 배치}
}

\author{
    이찬희 (Lee Chanhee)\\
    컴퓨터공학부\\
    서울과학기술대학교\\
    \texttt{happilyfly@seoultech.ac.kr}
}

\date{2026년 4월}

% ─── Document ────────────────────────────────────────────────────────────────
\begin{document}

\maketitle
\thispagestyle{empty}

% ─── Abstract ────────────────────────────────────────────────────────────────
\begin{abstract}
본 논문은 BTC/USDT 시간봉 시장에서 가격 방향 예측 없이 ATR 비례 지정가 주문 간격을
연속 액션으로 출력하는 \emph{동적 그리드 트레이딩} 강화학습 프레임워크를 제안한다.
기존 RL 트레이딩 연구의 주류인 방향성 전략 및 이산 액션 공간(매수/매도/홀딩)과 달리,
본 에이전트는 $(a_\text{agg}, a_\text{pt}) \in [0,1]^2$의 2차원 연속 액션을 통해
ATR(Average True Range)에 비례하는 4개의 지정가 주문 가격을 결정한다.
주문 배치 공식을 제어하는 8개의 계수는 베이지안 탐색(Optuna TPE, 50회 시도)으로
공동 최적화되어 기본값 대비 검증 Sharpe 비율 \(+145\%\) 향상을 달성했다.
2020--2022년 데이터로 300만 스텝 학습된 최종 모델은 검증 Sharpe \textbf{35.42},
테스트셋 Sharpe \textbf{43.04}(2024--2026)를 기록하며 최우수 고정 그리드 베이스라인을
\textbf{29.2배} 상회하는 동시에 최대낙폭 \textbf{3.1\%}를 유지했다.
단, 사후 행동 분석 결과 학습된 정책이 거의 상수 액션으로 수렴함이 확인되었으며,
이는 베이지안 최적화된 계수 자체가 성능의 주요 원동력임을 시사한다.
이 설계적 한계와 다음 단계 개선 방향을 함께 논의한다.
\end{abstract}

\vspace{1em}
\noindent\textbf{키워드:} 강화학습, 그리드 트레이딩, PPO, 베이지안 최적화,
암호화폐, ATR 비례 주문, 정책 포화

\newpage
\tableofcontents
\newpage

% ─── 1. Introduction ─────────────────────────────────────────────────────────
\section{서론}

암호화폐 시장은 장기 추세와 무관하게 반복적인 미세 가격 진동을 유발하는
높은 변동성을 지속적으로 보인다.
그리드 트레이딩은 이러한 특성을 활용하여 현재가 주변에 매수·매도 지정가 주문을
계단형으로 배치하고, 각 매수-매도 완료 사이클에서 수익을 획득한다
\citep{sun2023rl}.
기존 그리드 트레이딩의 핵심 약점은 배포 시점에 그리드 간격을 고정하고 이후 조정하지
않는다는 점이다: 간격이 너무 좁으면 수익성 없는 소규모 거래가 넘쳐나고,
너무 넓으면 대부분의 진동을 포착하지 못한다.

최근 RL 트레이딩 연구는 다음 가격이 오를지 내릴지 예측하는 \emph{방향성} 전략에
거의 전적으로 집중되어 있다 \citep{liu2021bitcoin, yasin2024rl, pham2025dqn}.
이 계열의 연구는 가격 예측의 어려움을 그대로 이어받으며,
예측 정확도가 낮은 횡보장이나 고변동성 시장에서 수익을 내기 어렵다.
RL 에이전트를 방향성 예측기가 아닌 \emph{시장 조성자} 또는 \emph{간격 선택자}로
다루는 연구는 극소수에 불과하여, 동적 그리드 문제는 대체로 미개척 영역으로 남아 있다.

본 논문은 다음 두 가지 연구 질문을 다룬다:

\begin{enumerate}[label=\textbf{RQ\arabic*:}]
    \item 연속 ATR 비례 액션 공간을 가진 PPO 에이전트가 미관측 BTC/USDT 데이터에서
          고정 간격 그리드 베이스라인을 능가할 수 있는가?
    \item 학습된 에이전트가 현재 변동성 레짐에 따라 그리드 간격을 적응적으로 조절하는가?
\end{enumerate}

\subsection*{기여 사항}

\begin{itemize}
    \item \textbf{비방향성 MDP 정식화.}
          시장 레짐과 포트폴리오 상태를 모두 인코딩하는 5차원 상태 공간에,
          방향성 가격 예측 없이 그리드 간격 폭을 제어하는 2차원 연속 액션을 결합했다.

    \item \textbf{ATR 비례 조정 가능 주문 배치 공식.}
          현재 ATR에 비례하는 거리에 4개의 지정가 주문(매수 2개, 매도 2개)을 배치한다.
          8개의 스칼라 계수가 공식을 매개변수화함으로써 경제적 해석을 유지하면서
          데이터 기반 보정을 가능하게 한다.

    \item \textbf{베이지안 계수 최적화.}
          Optuna TPE와 MedianPruner를 적용하여 8차원 계수 공간을 탐색,
          50회 시도만에 수작업 기본값 대비 검증 Sharpe $+145\%$ 향상을 달성했다.

    \item \textbf{엄격한 학습/검증/테스트 분리.}
          \citet{gort2022backtest}를 따라 테스트셋(2024--2026)은 모든 설계 결정이
          확정된 이후까지 완전히 봉인했으며, 테스트셋 평가는 정확히 한 번만 수행했다.
\end{itemize}

% ─── 2. Related Work ─────────────────────────────────────────────────────────
\section{관련 연구}

\subsection{RL 트레이딩 서베이}

\citet{sun2023rl}은 RL 트레이딩 문헌을 조사하여 두 가지 반복되는 한계를 지적한다:
학습된 정책의 해석 가능성 부족과 시장 레짐 간 낮은 일반화 성능.
\citet{liu2025review}는 2018--2025년에 걸친 10편의 논문을 리뷰하며,
RL 에이전트가 규칙 기반 방법 대비 비정상 시장에 대한 우수한 적응력을 보이지만
학습 안정성과 샘플 효율성이 실용적 배포의 주요 장벽으로 남아 있음을 확인한다.

\subsection{방향성 RL 트레이딩 시스템}

\citet{liu2021bitcoin}은 LSTM 가격 예측기와 PPO 에이전트를 결합하여 BTC/USDT에 적용하며,
예측된 가격을 상태 특징으로 입력한다.
이 2단계 파이프라인은 트레이딩 성능을 예측 정확도에 종속시켜,
예측 오류가 주문 결정에 직접 전파되는 문제를 안고 있다.
본 접근법은 예측 단계를 완전히 제거한다.

\citet{yasin2024rl}은 주식 시장에서 20개 기술 지표를 사용하여 DQN, PPO, A2C를
벤치마크한다.
이산 액션 공간(매수/매도/홀딩)과 관측에 포트폴리오 상태 정보가 없다는 점은,
현재 노출 대비 주문 크기를 조절해야 하는 그리드 트레이딩에서 치명적인 한계이다.

\citet{pham2025dqn}은 DQN을 \emph{전략 선택}에 적용한다: 에이전트가 전략 자체를
학습하는 대신 5개의 하드코딩된 규칙 기반 전략 중 하나를 매 스텝 선택한다.
지표 상태의 PCA 압축은 학습된 정책의 경제적 의미를 더욱 불분명하게 만든다.
본 연구에서는 에이전트가 그리드 간격을 직접 출력하도록 학습시켜 완전한
해석 가능성을 유지한다.

\subsection{위험 인식 RL 트레이딩}

\citet{bandarupalli2025risk}는 위험한 포지션을 억제하기 위해 PPO 보상에
명시적인 포트폴리오 분산 페널티를 추가한다.
이 페널티가 너무 많은 수익성 있는 거래를 억압하여 에이전트가 매수보유 전략보다
저조한 성과를 보인다고 보고한다.
이 발견은 별도의 위험 항을 두지 않고 자본 변화와 거래비용 페널티만으로 구성한
본 연구의 보상 설계를 뒷받침한다: 그리드 트레이딩은 본질적으로 변동성을 추구하므로
분산 페널티는 역효과를 낳는다.

\subsection{백테스팅 방법론}

\citet{gort2022backtest}는 RL 트레이딩 연구에서 단일 분할 검증이 상당한
백테스트 과적합으로 이어짐을 실증적으로 보인다.
이들의 엄격한 시간적 분할 권고가 본 연구의
학습(2020--2022) / 검증(2023) / 테스트(2024--2026) 설계에 직접 반영되었으며,
테스트 파티션은 모든 하이퍼파라미터 및 계수 결정이 완료된 후에만 접근했다.

\subsection{본 연구의 위치}

표~\ref{tab:comparison}은 가장 가까운 선행 연구와의 차이점을 요약한다.

\begin{table}[h]
\centering
\caption{관련 RL 트레이딩 시스템과의 비교.}
\label{tab:comparison}
\setlength{\tabcolsep}{5pt}
\begin{tabular}{lccccc}
\toprule
 & \textbf{본 연구} & Liu '21 & Yasin '24 & Pham '25 & Banda. '25 \\
\midrule
비방향성              & \checkmark & $\times$ & $\times$ & $\times$ & $\times$ \\
연속 액션 공간        & \checkmark & $\times$ & $\times$ & $\times$ & $\times$ \\
포지션 상태           & \checkmark & $\times$ & $\times$ & $\times$ & $\checkmark$ \\
ATR 비례 주문         & \checkmark & $\times$ & $\times$ & $\times$ & $\times$ \\
베이지안 계수 튜닝    & \checkmark & $\times$ & $\times$ & $\times$ & $\times$ \\
봉인된 테스트셋       & \checkmark & $\times$ & $\times$ & $\times$ & $\times$ \\
\bottomrule
\end{tabular}
\end{table}

% ─── 3. Methodology ──────────────────────────────────────────────────────────
\section{방법론}

\subsection{데이터 및 전처리}

CCXT 라이브러리를 통해 Binance에서 BTC/USDT 시간봉 캔들스틱 데이터를 수집하며,
2020년 1월부터 2026년 4월까지 중복 제거 후 54,733개 캔들이 포함된다.
데이터는 시간 순서에 따라 세 개의 비중첩 구간으로 분할된다:

\begin{itemize}
    \item \textbf{학습} (2020-01-01 -- 2022-12-31, 25,916캔들):
          상승장(2020--2021)과 주요 하락장(2022)을 모두 포함한다.
    \item \textbf{검증} (2023-01-01 -- 2023-12-31, 8,736캔들):
          하이퍼파라미터 선택 및 학습 중 모델 평가에 사용하는 회복기이다.
    \item \textbf{테스트} (2024-01-01 -- 2026-04-09, 19,901캔들):
          모든 설계 결정이 확정될 때까지 완전히 봉인 유지; 정확히 한 번만 평가한다.
\end{itemize}

각 원시 캔들에 롤링 윈도우 $W = 168$시간(1주) 기준으로 세 가지 파생 시계열을
추가한다:

\begin{align}
    \text{log\_price}_t &= \log\!\left(\frac{p_t}{\bar{p}_{t,W}}\right), \\
    \text{volatility}_t &= \frac{\text{ATR}(168)_t}{p_t}, \\
    \text{zscore}(x)_t  &= \frac{x_t - \mu_{x,t,W}}{\sigma_{x,t,W} + \epsilon},
\end{align}

여기서 $\bar{p}_{t,W}$는 롤링 평균 종가, ATR은 OHLC 데이터에서 직접 계산한
Average True Range, $\epsilon = 10^{-8}$은 수치 안정성을 보장한다.
정책 입력에서 척도 불일치를 방지하기 위해 \texttt{log\_price}와
\texttt{volatility}에 z-점수 변환을 적용하여 저장한다.

\subsection{트레이딩 환경 (MDP)}

트레이딩 문제를 Gymnasium 환경으로 구현된 마르코프 결정 과정
$\mathcal{M} = (\mathcal{S}, \mathcal{A}, \mathcal{T}, \mathcal{R}, \gamma)$으로 모델링한다.

\subsubsection*{상태 공간}

시각 $t$에서의 관측은 5차원 벡터이다:

\begin{equation}
    \bm{s}_t = \bigl[
        \underbrace{\hat{z}(\text{log\_price}_t)}_{\text{가격 수준}},\;
        \underbrace{d_t}_{\text{괴리도}},\;
        \underbrace{h_t}_{\text{포지션 크기}},\;
        \underbrace{c_t}_{\text{현금 비율}},\;
        \underbrace{\hat{z}(\text{volatility}_t)}_{\text{시장 변동성}}
    \bigr] \in \mathbb{R}^5,
\end{equation}

여기서 $\hat{z}(\cdot)$은 롤링 z-점수 정규화를 나타내며,

\begin{align}
    d_t &= \frac{\bar{p}^\text{avg}_t - p_t}{\bar{p}^\text{avg}_t}
           \quad \text{(평균 진입 가격 대비 괴리도)}, \\
    h_t &= \frac{q_t \cdot p_t}{K}
           \quad \text{(초기 자본 $K$ 대비 보유 가치 비율)}, \\
    c_t &= \frac{\text{cash}_t}{K}
           \quad \text{(초기 자본 대비 현금 비율)}.
\end{align}

$\bar{p}^\text{avg}_t$는 현재 사이클의 평균 진입 가격이며,
미보유 상태에서는 직전 완료 사이클의 평균 진입가를,
거래 이력이 없으면 0을 사용한다.
가격이 2배가 되면 동일한 BTC 보유량의 포트폴리오 위험도 2배가 되므로,
원시 수량 대신 가치 비율 $h_t$를 사용한다.

\subsubsection*{액션 공간}

액션 공간은 2차원 연속 박스이다:
\begin{equation}
    \mathcal{A} = [0, 1]^2,
    \quad
    \bm{a}_t = (a_\text{agg},\; a_\text{pt}),
\end{equation}

여기서 $a_\text{agg}$(공격성)는 매수 주문 배치를 제어하고
$a_\text{pt}$(수익 목표)는 매도 주문 배치를 제어한다.
이는 대부분의 선행 연구에서 사용하는 이산 액션 공간(매수/매도/홀딩)
\citep{liu2021bitcoin, yasin2024rl}과 극명히 대조되며,
이산 공간은 공격성의 \emph{정도}를 표현할 수 없다.

\subsubsection*{주문 배치 공식}

각 스텝에서 ATR 비례 공식을 통해 액션으로부터 4개의 지정가 주문이 계산된다:

\begin{align}
    r_t &= \frac{\text{ATR}(168)_t}{p_t}
           \quad \text{(정규화된 변동성 비율)}, \label{eq:atr_ratio} \\[4pt]
    \Delta^\text{buy-hi}_t  &= r_t \cdot (A_b + a_\text{agg}\cdot B_b), \label{eq:buy_hi} \\
    \Delta^\text{buy-lo}_t  &= r_t \cdot (C_b + a_\text{agg}\cdot D_b), \label{eq:buy_lo} \\
    \Delta^\text{sell-mkt}_t &= r_t \cdot (A_s + a_\text{pt}\cdot B_s), \label{eq:sell_mkt} \\
    \Delta^\text{sell-cst}_t &= r_t \cdot (C_s + a_\text{pt}\cdot D_s), \label{eq:sell_cst}
\end{align}

그리고 4개의 주문 가격은:

\begin{align}
    P^\text{buy-hi}_t  &= p_t\cdot(1 - \Delta^\text{buy-hi}_t), \\
    P^\text{buy-lo}_t  &= p_t\cdot(1 - \Delta^\text{buy-lo}_t), \\
    P^\text{sell-mkt}_t &= p_t\cdot(1 + \Delta^\text{sell-mkt}_t), \\
    P^\text{sell-cst}_t &= \bar{p}^\text{avg}_t\cdot(1 + \Delta^\text{sell-cst}_t).
\end{align}

$\bm{\theta} = (A_b, B_b, C_b, D_b, A_s, B_s, C_s, D_s) \in \mathbb{R}^8$은
공식 계수이다.
ATR 비례 스케일링은 현재 변동성 수준에 무관하게 액션 범위 $[0,1]$이
의미 있는 체결 확률 스펙트럼에 일관되게 대응하도록 보장하여,
대부분의 지정가 주문이 체결되지 않는 레짐을 방지한다.

이중 매도 주문 설계는 두 가지 보완적 청산 채널을 반영한다:
\texttt{sell-mkt}는 현재가 위에서 빠른 청산을 목표로 하여 신속한 현금 회수를 달성하고,
\texttt{sell-cst}는 평균 원가 이상에서 큰 이익을 목표로 하여 대규모 랠리를 포착한다.

\subsubsection*{주문 체결 및 사이클 로직}

주문 체결은 다음 캔들의 고가/저가(차봉 체결 모델)로 시뮬레이션되며,
같은 봉에서는 매도 주문을 매수 주문보다 먼저 처리한다.
각 매수-매도 완료 시퀀스가 하나의 \emph{트레이딩 사이클}을 구성한다.
사이클 내 자본 배분은 고정 예산 방식을 사용한다: 사이클 시작 시 가용 현금을
$n_\text{splits} = 2$ 슬롯으로 나누고, 각 슬롯이 $n_\text{buy} = 2$개의
지정가 주문에 자금을 지원한다.
총 보유량이 단일 슬롯 가치 이하로 떨어지면 남은 보유량을 전량 청산하여
사이클 종료를 보장한다.

\subsubsection*{보상}

각 스텝의 보상은 정규화된 자본 변화량이다:

\begin{equation}
    r_t = \frac{V_t - V_{t-1}}{K},
    \label{eq:reward}
\end{equation}

여기서 $V_t = \text{cash}_t + q_t \cdot p_t$는 총 포트폴리오 가치이다.
거래비용(건당 0.05\%, Binance maker 수수료)은 체결 시점에 현금 잔액에서 직접
차감되므로, 별도의 페널티 항 없이 자본 변화에 자동으로 반영된다.

\subsection{PPO 에이전트}

Stable-Baselines3를 통해 Proximal Policy Optimization \citep{schulman2017ppo}을 사용한다.
정책 네트워크와 가치 네트워크는 MLP 트렁크(64유닛 은닉층 2개, tanh 활성화)를 공유한다.
액션 분포는 sigmoid 스쿼싱 레이어로 $[0, 1]^2$에 클리핑된 대각 가우시안이다.

학습은 4개의 병렬 환경(\texttt{n\_envs}~$= 4$)에서 에피소드 길이
$T_\text{ep} = 2016$ 스텝($\approx 12$주)으로 진행된다.
학습 윈도우 내 무작위 에피소드 시작이 에이전트를 다양한 시장 레짐에 노출시킨다.
코사인 학습률 스케줄이 300만 환경 스텝에 걸쳐 $3\times10^{-4}$에서 $10^{-5}$로
감소한다.
엔트로피 계수 $\alpha_H = 0.05$는 절제 실험으로 선택되었다
(절 \ref{sec:param-search}).

\subsection{베이지안 계수 최적화}
\label{sec:bayesian}

$\bm{\theta}$의 8개 공식 계수는 사전 경제적 지식을 인코딩하지만
(예: $A_b < C_b$는 buy-hi가 buy-lo보다 현재가에 가깝도록 강제),
정확한 값은 알 수 없다.
이를 하이퍼파라미터로 취급하여 MedianPruner를 갖춘 Tree-structured Parzen Estimator
(TPE)로 8차원 연속 공간을 탐색한다:

\begin{itemize}
    \item \textbf{목적 함수}: 워밍업된 PPO 에이전트로 $10^6$ 학습 스텝 후
          5개 에피소드에 대한 평균 검증 Sharpe 비율.
    \item \textbf{소프트 제약}: $C_b > A_b$, $D_b > B_b$ 위반 시도(buy-hi/buy-lo
          순서 역전) 시 센티넬 값 $-999$를 반환하고 TPE 모델에서 제외.
    \item \textbf{가지치기}: 워밍업 10회 후 MedianPruner가 비전망적 시도를
          첫 중간 평가 후 조기 종료하여 유망한 구성에 계산 자원 회수.
    \item \textbf{예산}: 50회 시도($\approx$ 단일 CPU 4.3시간).
\end{itemize}

탐색 공간과 최종 최적화 값은 표~\ref{tab:coefs}에 제시한다.

\begin{table}[h]
\centering
\caption{공식 계수 탐색 공간 및 베이지안 최적화 값 (Trial \#42).}
\label{tab:coefs}
\setlength{\tabcolsep}{6pt}
\begin{tabular}{llrrrr}
\toprule
계수 & 역할 & 탐색 최솟값 & 탐색 최댓값 & \textbf{기본값} & \textbf{최적값} \\
\midrule
$A_b$ & buy-hi 기저 & 0.1 & 3.0 & 0.500 & \textbf{0.285} \\
$B_b$ & buy-hi 범위 & 0.1 & 5.0 & 1.500 & \textbf{1.748} \\
$C_b$ & buy-lo 기저 & 0.1 & 10.0 & 2.500 & \textbf{5.223} \\
$D_b$ & buy-lo 범위 & 0.5 & 20.0 & 7.500 & \textbf{18.683} \\
$A_s$ & sell-mkt 기저 & 0.1 & 3.0 & 0.500 & \textbf{0.101} \\
$B_s$ & sell-mkt 범위 & 0.1 & 5.0 & 1.500 & \textbf{0.890} \\
$C_s$ & sell-cost 기저 & 0.1 & 10.0 & 2.500 & \textbf{6.913} \\
$D_s$ & sell-cost 범위 & 0.5 & 20.0 & 7.500 & \textbf{5.457} \\
\bottomrule
\end{tabular}
\end{table}

% ─── 4. Experimental Setup ───────────────────────────────────────────────────
\section{실험 설정}

\subsection{베이스라인}

학습을 사용하지 않는 7개의 고정 규칙 전략과 비교한다:

\begin{itemize}
    \item \textbf{매수보유(Buy-and-Hold)}: 시작 시 BTC 매수, 종료 시 매도.
    \item \textbf{고정 그리드 $\delta\%$} ($\delta \in \{1, 2, 5\}$):
          현재가 $\delta\%$ 아래에 매수 주문, $\delta\%$ 위에 매도 주문 배치.
    \item \textbf{ATR 그리드 $k$} ($k \in \{0.5, 1.0, 2.0\}$):
          주문 간격이 현재 ATR의 $k$배로, 에이전트의 ATR 비례 스케일링을
          학습 없이 모방한다.
\end{itemize}

ATR 그리드 베이스라인은 특히 유익한 비교 대상이다: 에이전트와 동일한 시장 정보(ATR)를
공유하지만 고정 규칙으로 적용한다.
이 베이스라인 대비 성능 차이는 베이지안 계수 최적화와 RL 학습 절차의 복합적 효과를
반영한다.

\subsection{평가 지표}

주요 지표는 연율화 Sharpe 비율이다:

\begin{equation}
    \text{Sharpe} = \frac{\mathbb{E}[r_t]}{\text{Std}(r_t)} \cdot \sqrt{8760},
    \label{eq:sharpe}
\end{equation}

여기서 $r_t$는 스텝별 수익률이고 8760은 연간 시간봉 수이다.
보조 지표로 최대낙폭(MDD), 총 수익률, 완료된 트레이딩 사이클 수,
사이클 승률을 포함한다.

\subsection{실험 프로토콜}

모든 실험에 시드 42를 사용한다.
검증 Sharpe는 검증 윈도우 내 무작위 시작으로 길이 $T_\text{ep} = 2016$ 스텝의
5개 평가 에피소드에 대한 평균으로 계산되며, 결정론적($\text{deterministic}=\text{True}$)
추론을 사용한다.
테스트 평가는 전체 19,901 스텝 테스트 시퀀스에 대한 단일 결정론적 패스로,
무작위 시작 선택으로부터 어떤 정보도 보고 결과에 영향을 미칠 수 없음을 보장한다.

% ─── 5. Results ──────────────────────────────────────────────────────────────
\section{실험 결과}

\subsection{하이퍼파라미터 탐색}
\label{sec:param-search}

베이지안 계수 튜닝 전에 검증셋에서 세 가지 환경 하이퍼파라미터에 대한
체계적인 절제 실험을 수행한다:

\begin{itemize}
    \item \textbf{자본 슬롯 수} $n_\text{splits} \in \{2, 8, 16\}$:
          $n_\text{splits}=2$가 사이클당 자본을 집중하고 유휴 현금을 줄여 가장 높은
          Sharpe를 달성한다.
    \item \textbf{매도 임계 기준} (\texttt{avg\_price} vs.\ \texttt{price}):
          매도 비용 주문을 현재 시장가가 아닌 사이클 평균 진입가에 고정하면
          더 일관적인 수익 확정이 가능하다.
    \item \textbf{엔트로피 계수} $\alpha_H \in \{0.01, 0.05, 0.10\}$:
          $\alpha_H = 0.05$가 탐색과 활용의 균형을 맞추며,
          Sharpe를 14.39(VecNorm 없는 기준, exp007)에서 17.58로 향상시킨다.
\end{itemize}

기본 공식 계수로 선택된 구성(exp013b)은 검증 Sharpe 17.58을 달성하며,
이것이 베이지안 최적화의 기준값으로 사용된다.

\subsection{베이지안 계수 최적화}

그림~\ref{fig:optuna}은 50회 베이지안 탐색 결과를 보여준다.
최우수 시도(\#42)는 검증 Sharpe 42.997을 달성했으며, 이는 기본 계수 기준값 대비
$+145\%$ 향상이다.

\begin{figure}[H]
    \centering
    \begin{subfigure}[b]{0.48\textwidth}
        \centering
        \includegraphics[width=\textwidth]{fig_01_distribution.png}
        \caption{Sharpe distribution across 50 trials.
                 Red dashed line: default-coefficient baseline (17.58).}
        \label{fig:optuna_dist}
    \end{subfigure}
    \hfill
    \begin{subfigure}[b]{0.48\textwidth}
        \centering
        \includegraphics[width=\textwidth]{fig_03_importance.png}
        \caption{Pearson correlation between each coefficient and Sharpe ratio.
                 $A_s$ (sell-market base) shows the strongest negative correlation.}
        \label{fig:optuna_importance}
    \end{subfigure}
    \caption{Bayesian coefficient tuning results (50 TPE trials, 1M steps each).}
    \label{fig:optuna}
\end{figure}

계수 중요도 분석(그림~\ref{fig:optuna_importance})은 매도 시장가 주문 간격의 기저값인
$A_s$가 단연 가장 영향력 있는 파라미터임을 보여준다.
최적값(0.101)이 탐색 하한(0.1)에 위치한다는 점은,
에이전트가 매우 타이트한 매도 시장가 주문, 즉 현재가 위 $0.101 \times r_t$ 거리에
빠른 청산 호가를 배치하여 각 매수 체결 후 거의 즉각적인 현금 회수를 가능하게 함으로써
이익을 얻음을 시사한다.

그림~\ref{fig:gap_comparison}은 최적화된 계수가 기본값 대비 효과적인 주문 간격
범위를 어떻게 변화시키는지 보여준다.

\begin{figure}[H]
    \centering
    \includegraphics[width=0.85\textwidth]{fig_06_gap_comparison.png}
    \caption{Effective order-gap ranges (in ATR multiples) for default vs.\ optimized
             coefficients, shown across the full action range $[0,1]$.
             Note the dramatic expansion of the buy-lo range and the sharp contraction
             of the sell-market gap.}
    \label{fig:gap_comparison}
\end{figure}

최적화된 전략은 \emph{빠른 매도 / 깊은 매수} 정책으로 해석할 수 있다:
\begin{itemize}
    \item \textbf{sell-mkt} 간격이 $[0.10, 1.0] \times r_t$로 축소
          (기본값 $[0.5, 2.0] \times r_t$ 대비): 소폭 상승 시에도 즉시 청산하여
          사이클 빈도를 극대화한다.
    \item \textbf{buy-lo} 간격이 $[5.2, 23.9] \times r_t$로 확장
          (기본값 $[2.5, 10.0] \times r_t$ 대비): 고변동성 이벤트에서 나타나는
          큰 하락을 포착하기 위해 깊은 대기 매수를 배치한다.
    \item \textbf{sell-cost} 간격이 $[6.9, 12.4] \times r_t$로 상승:
          원가 기준 매도는 평균 진입가 이상의 대규모 랠리에서만 발동된다.
\end{itemize}

\subsection{최종 모델 학습 (exp016)}

Trial \#42의 최적화된 계수를 사용하여 최종 PPO 에이전트를 300만 스텝 학습시킨다.
그림~\ref{fig:learning_curve}은 학습 전 과정의 검증 Sharpe를 보여준다.

\begin{figure}[H]
    \centering
    \includegraphics[width=\textwidth]{03_learning_curve.png}
    \caption{Validation Sharpe ratio vs.\ training steps.
             \textbf{(a) Pre-fix experiments} (exp001--002): a reward-function bug
             (fee double-counting) and an eval-pipeline bug kept Sharpe below 1.0.
             \textbf{(b) Post-fix experiments}: after both bugs were resolved, Sharpe
             climbed from $\sim$12 (exp008) to 17.6 (exp013b) through hyperparameter
             search, and finally to 35.4 with Bayesian-optimized coefficients (exp016,
             best checkpoint at step 2.1M, $\star$).}
    \label{fig:learning_curve}
\end{figure}

최선 체크포인트는 스텝 2,100,000에서 검증 Sharpe 35.42, MDD 2.46\%로 저장된다.
전체 학습 Sharpe(35.42)가 1M 스텝 Optuna 목적값(42.997)보다 낮은 이유는,
Optuna 평가는 5개의 무작위 시작 에피소드를 사용하여 유리한 시장 윈도우를 선택할 수
있는 반면, 최종 평가는 고정된 평가 윈도우에 걸쳐 평균을 취하여 더 보수적이고
재현 가능한 추정치를 제공하기 때문이다.

\subsection{검증셋 결과}

표~\ref{tab:val_results}는 2023년 검증셋에서 최선 PPO 체크포인트와 모든 베이스라인을
비교한다.

\begin{table}[H]
\centering
\caption{검증셋 결과 (2023년, 8,736캔들).}
\label{tab:val_results}
\setlength{\tabcolsep}{6pt}
\begin{tabular}{lrrrr}
\toprule
전략 & Sharpe & 수익률 (\%) & MDD (\%) & 사이클 수 \\
\midrule
매수보유              &  2.377 & +150.18 & 21.74 &    0 \\
고정 그리드 1\%       &  1.991 &  +64.11 & 17.50 &  120 \\
고정 그리드 2\%       &  2.039 &  +31.05 & 12.57 &   38 \\
고정 그리드 5\%       &  1.377 &   +4.99 &  3.32 &    4 \\
ATR 그리드 $k{=}0.5$ &  1.141 &  +30.80 & 20.66 &  279 \\
ATR 그리드 $k{=}1.0$ &  1.607 &  +54.37 & 18.63 &  177 \\
ATR 그리드 $k{=}2.0$ &  2.234 &  +55.42 & 13.23 &   91 \\
\midrule
\textbf{PPO (exp016)} & \textbf{35.424} & \textbf{+365.60} & \textbf{2.46} & \textbf{938} \\
\bottomrule
\end{tabular}
\end{table}

PPO 에이전트는 최강 베이스라인(ATR 그리드 $k{=}2.0$, Sharpe 2.234) 대비
MDD를 13.23\%에서 2.46\%로 낮추면서 \textbf{15.9배} 향상된 Sharpe를 달성했다.

\subsection{테스트셋 결과}
\label{sec:test-results}

테스트셋(2024-01-01 -- 2026-04-09)은 모든 모델 및 계수 결정이 확정된 후
단 한 번 접근했다.
이 기간은 BTC 가격이 약 \$42,000에서 \$100,000 이상으로 상승 후 \$75,000로
하락하는 과정을 포함하며, 2024년 주요 강세장과 2025년 1분기 급격한 조정이
모두 포함된 다양한 표본 외 평가 구간이다.

\begin{table}[H]
\centering
\caption{테스트셋 결과 (2024--2026년, 19,901캔들).}
\label{tab:test_results}
\setlength{\tabcolsep}{6pt}
\begin{tabular}{lrrrr}
\toprule
전략 & Sharpe & 수익률 (\%) & MDD (\%) & 사이클 수 \\
\midrule
매수보유              &  0.681 &  +62.25 & 50.08 &    0 \\
고정 그리드 1\%       &  0.724 &  +52.81 & 29.87 &  176 \\
고정 그리드 2\%       &  0.625 &  +34.76 & 26.66 &   54 \\
고정 그리드 5\%       &  1.472 &  +14.28 &  3.16 &    9 \\
ATR 그리드 $k{=}0.5$ & $-$0.004 & $-$20.22 & 49.77 &  325 \\
ATR 그리드 $k{=}1.0$ &  0.258 &   +5.29 & 49.59 &  242 \\
ATR 그리드 $k{=}2.0$ &  0.436 &  +22.73 & 29.85 &   88 \\
\midrule
\textbf{PPO (exp016)} & \textbf{43.040} & --- & \textbf{3.12} & \textbf{9,520} \\
\bottomrule
\end{tabular}
\end{table}

\noindent\textit{테스트셋 수익률 주석: 9,520 완료 사이클, 평균 사이클 PnL 0.182\%가
27개월에 걸쳐 복리로 누적되면 산술 누적 수익률이 천문학적 값에 도달한다.
시간과 위험을 정규화하는 Sharpe 비율과 MDD가 적절한 지표이며,
3.12\% MDD가 자본 안전성의 실질적인 척도이다.}

그림~\ref{fig:equity}은 모든 베이스라인 전략과 함께 로그 스케일로 테스트 기간
자본 곡선을 보여준다.

\begin{figure}[H]
    \centering
    \includegraphics[width=0.9\textwidth]{06_test_equity_curve.png}
    \caption{Portfolio equity curves on the test set (2024--2026), log scale.
             The PPO agent (blue) grows exponentially while all baselines
             remain near the initial capital line. Bottom panel: BTC/USDT price.}
    \label{fig:equity}
\end{figure}

그림~\ref{fig:progression}은 전체 실험 진행 맥락에서 테스트 결과를 위치시킨다.

\begin{figure}[H]
    \centering
    \includegraphics[width=0.9\textwidth]{06_sharpe_progression.png}
    \caption{\textbf{Left}: validation Sharpe at each major experiment milestone.
             \textbf{Right}: final PPO Sharpe vs.\ best baseline on Val and Test sets.
             Test Sharpe exceeds Val Sharpe, providing no evidence of overfitting.}
    \label{fig:progression}
\end{figure}

% ─── 6. Analysis ─────────────────────────────────────────────────────────────
\section{분석 및 논의}

\subsection{RQ1: PPO가 고정 그리드 베이스라인을 능가하는가?}

검증셋에서 PPO는 Sharpe 35.42를 달성하여 최우수 베이스라인 2.23 대비
\textbf{15.9배} 우위를 보인다.
봉인 테스트셋에서 PPO는 43.04를 달성하여 고정 그리드 5\% 대비 \textbf{29.2배}의
우위를 보인다.
두 결과 모두 RQ1을 명확히 긍정적으로 답한다.

검증 Sharpe 대비 높은 테스트 Sharpe(43.04 vs.\ 35.42)는 2024--2025년 BTC 강세장이
에이전트의 고빈도 빠른 사이클 트레이딩 스타일에 특히 유리한 환경을 제공했기 때문이다.
핵심적으로, 이 차이의 방향(테스트 $>$ 검증)은 테스트셋 과적합을 배제한다:
테스트 데이터를 암묵적으로 본 모델은 검증 대비 \emph{더 나을 수 없는} 테스트 성능을
보일 것이기 때문이다.

\subsection{RQ2: 변동성 레짐 적응}

그림~\ref{fig:behavior}은 ATR/가격 임계값으로 정의된 변동성 레짐(낮음/중간/높음)에
따라 계층화된 테스트셋에서의 에이전트 액션 출력을 보여준다.

\begin{figure}[H]
    \centering
    \includegraphics[width=\textwidth]{06_behavior_analysis.png}
    \caption{Action distributions on the test set.
             \textbf{Top-left/Bottom-left}: both aggressiveness and profit target
             are concentrated at exactly 0.0 across all volatility regimes,
             indicating a saturated constant policy.
             \textbf{Bottom-right}: the regime distribution across the test period.}
    \label{fig:behavior}
\end{figure}

RQ2의 기대와 달리, 학습된 정책은 관측된 변동성 레짐에 무관하게 결정론적으로
거의 상수인 액션 벡터 $\bm{a} \approx (0, 0)$을 출력한다.
정책 네트워크 검사로 이를 확인한다: 공격성의 원시 가우시안 평균은 $-9.19$
(tanh 스쿼싱 후 $\approx -1$, $[0,1]$ 재스케일링 후 $\approx 0$)이고,
수익 목표는 $-4.30$(재스케일링 후 $\approx 0$)이다.
낮음, 중간, 높음 변동성 레짐에 걸쳐 통계적으로 유의미한 액션 값 차이는 존재하지 않는다.

\textbf{왜 정책이 포화되었는가?}
근본 원인은 베이지안 계수 최적화와 RL 액션 공간 사이의 구조적 결합에 있다.
주문 배치 공식은:
\begin{equation}
    \Delta^\text{sell-mkt}_t = r_t \cdot (A_s + a_\text{pt} \cdot B_s).
\end{equation}
베이지안 탐색이 $A_s$를 하한값(0.101)으로 이끌었으며, 이는 최솟값 액션
($a_\text{pt} = 0$)을 이미 최적으로 만들었다: $a_\text{pt}$를 증가시키면
매도 간격만 넓어져 사이클 빈도가 감소할 뿐이다.
RL 에이전트는 $a_\text{pt} = 0$을 유지하는 것이 올바름을 학습했지만,
이는 상태 기반 적응을 불필요하게 만든다.
달리 말하면, 베이지안 단계가 계수 수준에서 간격 결정 문제를 해결하여
RL이 활용할 잔여 자유도가 남지 않은 것이다.

\textbf{함의.}
현재 설계로는 RQ2에 답할 수 없다.
이 결과가 변동성 레짐 적응이 불가능하거나 도움이 안 된다는 의미가 아니라,
계수 최적화가 적응이 표현될 메커니즘을 제거했음을 의미한다.
계수가 행동의 \emph{범위}를 고정하고 RL이 \emph{그 범위 내 선택}을 제어하는
재설계된 시스템이 자연스러운 다음 단계이다
(절 \ref{sec:conclusion}).

\subsection{사이클 행동}

그림~\ref{fig:cycle}은 테스트 기간 동안 완료된 9,520개 트레이딩 사이클을 특성화한다.

\begin{figure}[H]
    \centering
    \includegraphics[width=\textwidth]{06_cycle_analysis.png}
    \caption{Test-set cycle analysis.
             \textbf{Left}: PnL distribution (win rate 96.9\%, mean +0.182\%).
             \textbf{Center}: cycle duration distribution (mean 1.2 bars $\approx 1.2$ hours).
             \textbf{Right}: cumulative compounding of per-cycle returns on a log scale.}
    \label{fig:cycle}
\end{figure}

96.9\% 사이클 승률과 평균 1.2시간 지속 시간은 에이전트가 빠르게 포지션을 열고 닫으며
각 사이클에서 소액 스프레드를 수집하는 준시장조성 모드로 운영됨을 보여준다.
오른쪽 패널의 로그-선형 누적 수익 곡선은 3.12\% MDD와 일관되게 급격한 낙폭 없이
복리가 꾸준히 발생함을 확인한다.

\subsection{한계}

\textbf{시장 충격.}
시뮬레이션은 무한 유동성을 가정한다: 지정가 주문은 명시된 가격에 슬리피지 없이
체결된다.
시뮬레이션 자본 규모(\$10,000)에서 이 가정은 BTC/USDT(일 거래량 $>$\$300억)에
합리적이다.
그러나 더 큰 포지션으로 확장하면 시장 충격 모델이 필요하다.

\textbf{수수료 모델.}
고정 0.05\% maker 수수료가 적용된다.
실제로는 taker 수수료(0.10\%), 펀딩 비용, 스프레드 비용이 특히 많은 체결을
발생시키는 빠른 사이클 전략에서 수익을 추가로 잠식할 것이다.

\textbf{학습 기간.}
학습셋(2020--2022)은 하나의 주요 강세-약세 사이클만 포함한다.
추가적인 시장 레짐의 확장된 커버리지가 일반화 주장을 강화할 것이다.

\textbf{정책 포화.}
학습된 정책이 거의 상수인 $(0, 0)$ 액션을 출력하는 이유는, 베이지안 최적화가
$A_s$를 하한으로 이끌어 최솟값 액션을 영구적으로 최적화하기 때문이다.
이는 구조적 설계 한계이다: 공식 계수와 RL 정책을 동일한 목적으로 공동 최적화하면
정책의 자유도가 제거된다.
향후 연구는 이 두 최적화 단계를 분리해야 한다.

\textbf{$A_s$ 경계.}
최적 $A_s = 0.101$이 베이지안 탐색 공간의 하한에 위치한다는 것은,
베이지안 단계가 매도 간격 문제를 완전히 해결하여 RL이 추가 개선에 기여할
여지를 남기지 않았음을 시사한다.

% ─── 7. Conclusion ───────────────────────────────────────────────────────────
\section{결론}
\label{sec:conclusion}

본 논문은 방향성 가격 예측 없이 PPO 에이전트가 ATR 비례 지정가 주문 배치를 제어하는
BTC/USDT 동적 그리드 트레이딩을 위한 비방향성 강화학습 프레임워크를 제안했다.
주요 기여는: (1) 시장 조건과 포트폴리오 상태를 모두 인코딩하는 5차원 MDP 정식화;
(2) 주문 의미론의 해석 가능성을 유지하는 8계수 ATR 스케일링 공식;
(3) 수작업 기본값 대비 검증 Sharpe를 $+145\%$ 향상시킨 베이지안 탐색 절차이다.

최종 모델은 테스트셋 Sharpe 43.04 — 최우수 고정 그리드 베이스라인의 29.2배 —
와 3.12\% 최대낙폭을 달성했다.
그러나 사후 분석에서 정책이 거의 상수인 $(0, 0)$ 액션 출력으로 수렴함이 밝혀졌다:
베이지안 최적화가 매도 시장가 기저 계수 $A_s$를 하한으로 이끌어 최솟값 액션을
영구적으로 최적화하고 상태 기반 선택의 필요성을 제거했다.
따라서 고정 베이스라인 대비 성능 우위는 학습된 적응적 행동보다 주로
베이지안 최적화 계수 집합에서 비롯된다.

이 발견은 연구 기여를 재구성한다: RL이 시장 레짐에 \emph{적응}함을 보이는 대신,
본 연구는 (a) 베이지안 최적화 계수를 가진 ATR 비례 프레임워크가 그 자체로
강력한 고정 전략을 구성하며, (b) 계수와 RL 정책의 공동 최적화가 레짐 적응의
출현을 막는 구조적 충돌을 일으킴을 보인다.

\textbf{다음 단계 설계.}
진정한 레짐 적응을 가능하게 하려면 세 가지 변경이 필요하다:
\begin{enumerate}
    \item \textbf{계수와 액션 역할 분리.}
          계수는 경제적으로 의미 있는 \emph{범위}를 고정해야 하며
          (예: 매도 시장가 간격을 $[0.1, 1.0]\times r_t$ 범위로 고정),
          베이지안 결과를 최적화 값이 아닌 정보에 입각한 경계로 활용한다.
          RL 액션은 그 후 상태에 기반하여 해당 범위 \emph{내에서} 선택한다.
    \item \textbf{방향성 상태 특징 추가.}
          정책이 강세, 약세, 횡보 레짐을 구별하는 데 필요한 신호를 갖도록
          롤링 로그 수익률 특징(예: 1일 및 1주일 추세)을 포함한다.
    \item \textbf{학습 중 정책 엔트로피 모니터링.}
          각 평가 체크포인트에서 액션 표준편차를 추적하며,
          0으로의 수렴은 여기서 식별된 포화 실패 모드의 조기 경보이다.
\end{enumerate}

추가 확장으로는 연속 액션 공간에서의 향상된 샘플 효율을 위한 SAC 도입,
그리고 실제 시장 조건(네트워크 지연 및 부분 체결 포함)에서의 검증을 위한
Binance 실시간 API와의 통합이 있다.

% ─── References ──────────────────────────────────────────────────────────────
\newpage
\bibliographystyle{plainnat}
\bibliography{references}

\end{document}
